\documentclass[12pt]{article}
\usepackage[T1]{fontenc}
\usepackage[utf8]{inputenc}
\usepackage{lmodern}
\usepackage{textcomp}
\usepackage{enumitem}
\usepackage{geometry}
\usepackage{graphicx}
\usepackage{amsmath, amssymb}
\geometry{margin=1in}
\begin{document}
\begin{center}
Mathematics - Graduate Level Assessment 20 \
Total Marks: 40 \
Answer all questions.
\end{center}

\section*{Multiple Choice (10 marks)}
\begin{enumerate}[label=\arabic*.]
\item If $m$ and $n$ are odd integers, how many terms in the expansion of $(m+n)^6$ are odd?
\begin{enumerate}[label=(\alph*)]
    \item 6
    \item 3
    \item All the terms are odd
    \item 4
    \item 1
\end{enumerate}
\item If f is a linear function with nonzero slope, and c \textless{} d, which of the following must be FALSE?
\begin{enumerate}[label=(\alph*)]
    \item f (d) \textgreater{} f (c)
    \item f (c) \textless{} f (d)
    \item f (d) \ensuremath{\neq} f (c)
    \item f (c) = f (d)
    \item f (d) \textless{} f (c)
\end{enumerate}
\item Statement 1 | Every solvable group is of prime-power order. Statement 2 | Every group of prime-power order is solvable.
\begin{enumerate}[label=(\alph*)]
    \item Statement 1 is always true, Statement 2 is false only for certain groups
    \item True, True
    \item Statement 1 is always false, Statement 2 is true only for certain groups
    \item True, False
    \item Statement 1 is true only for certain groups, Statement 2 is always true
\end{enumerate}
\item Solve the equation -6 x = -36. Check your answer.
\begin{enumerate}[label=(\alph*)]
    \item 42
    \item 30
    \item -42
    \item -36
    \item 1 over 216
\end{enumerate}
\item What is the order of group Z\_\{18\}?
\begin{enumerate}[label=(\alph*)]
    \item 30
    \item 15
    \item 12
    \item 6
    \item 24
\end{enumerate}
\end{enumerate}

\section*{True / False (10 marks)}
\textit{Answer True or False}
\begin{enumerate}[label=\arabic*.]
\setcounter{enumi}{5}
\item True or False: Calculating 60\% of 30 yields the result of 18, which is a common mistake in percentage calculations.
\item True or False: The element c = 2 in Z\_3 ensures that Z\_3[x]/($x^3$ + $x^2$ + c) is a field.
\item True or False: The polynomial 8 $x^3$ + 6 $x^2$ - 9 x + 24 is irreducible over Z but does not satisfy the Eisenstein criterion.
\item True or False: The covariance function C\_N(2, 4) for a Poisson process with rate λ = 4 is 12.
\item True or False: The solution to the equation v - 26 = 68 is v = 94.
\end{enumerate}

\section*{Long Form Response (20 marks)}
\textit{Answer each question with a clear and complete explanation.}
\begin{enumerate}[label=\arabic*.]
\setcounter{enumi}{10}
\item Find the roots of \[6 x^4 - 35 x^3 + 62 x^2 - 35 x + 6 = 0.\]Enter the roots, separated by commas.
\item Discuss how the amount of \$20.08 can be represented as a fraction, and analyze the implications of this representation in mathematical contexts.
\end{enumerate}

\vfill
\noindent\textit{End of Paper}
\end{document}